%%%%%%%%%%  scan idx 611  |  folio 598  %%%%%%%%%%
Donc \(P\to P/R\), qui fait de \(P\) un \(R\)-torseur sur \(P/R\), est affine. D'ailleurs on sait (SGA 3 1.4.2) que \(P/R\to Q\) est une immersion fermée, donc affine, donc il en est de même du composé \(P\to P/R\to Q\).

\textbf{Corollaire 1.5.} \uline{Soit} \(k\) \uline{un corps.}
\begin{enumerate}[label=\alph*)]
\item \uline{Soit} \(X\) \uline{un} \(k\)\uline{-schéma propre. Alors} \(\underline{\Pic}_{X/k}\) \uline{est représentable} (J. P. Murre).
\item \uline{Soit} \(f:X\to Y\) \uline{un morphisme surjectif entre deux} \(k\)\uline{-schémas propres. Alors le morphisme induit} \(f^*:\underline{\Pic}_{Y/k}\longrightarrow \underline{\Pic}_{X/k}\) \uline{est affine.}
\end{enumerate}

\textbf{Corollaire 1.6.} \uline{Soient} \(S\) \uline{un schéma noethérien intègre,} \(f:X\to Y\) \uline{un} \(S\)\uline{-morphisme surjectif entre deux} \(S\)\uline{-schémas propres.}
\begin{enumerate}[label=\alph*)]
\item \uline{Si} \(Y\) \uline{est normal sur} \(S\) \uline{(c'est-à-dire plat à fibres géométriquement normales), alors, quitte à restreindre} \(S\), \(f^*:\underline{\Pic}_{Y/S}\to\underline{\Pic}_{X/S}\) \uline{est quasi-fini.}
\item \uline{Si} \(\underline{\Pic}_{Y/S}\) \uline{est essentiellement propre sur} \(S\) \uline{(c'est-à-dire, s'il satisfait au critère valuatif de propreté, par exemple, si} \(Y\) \uline{est lisse sur} \(S\)\uline{) alors quitte à restreindre} \(S\), \(f^*:\underline{\Pic}_{Y/S}\to\underline{\Pic}_{X/S}\) \uline{est fini (et en particulier, affine).}
\end{enumerate}

En effet dans a), pour tout \(s\in S\), \(R_s=\Ker(f_s^*)\) est un groupe algébrique affine~; sa composante neutre est un groupe algébrique propre sur \(k\) (FGA 236, Théorème 2.1 (iii))~; \(R_s\) est donc fini. Dans b) \(f^*\) est à la fois propre et quasi-affine, donc fini.

\textbf{Remarque 1.7.} On ignore si 1.1 reste valable quand on y remplace « quasi-affine » par affine, comme pourraient le suggérer 1.5 et 1.6.

